\documentclass[10pt,twocolumn]{article}

\usepackage[letterpaper,margin=0.75in]{geometry}
\usepackage{times}
\usepackage{microtype}
\usepackage{booktabs}
\usepackage{array}
\usepackage{xcolor}
\usepackage{listings}
\usepackage{titlesec}
\usepackage{enumitem}
\usepackage{url}
\definecolor{linkblue}{rgb}{0.0,0.18,0.55}
\usepackage{hyperref}
\hypersetup{colorlinks=true,urlcolor=linkblue,linkcolor=black,citecolor=black,breaklinks=true}
\usepackage{caption}
\usepackage{abstract}

\captionsetup{font=small,labelfont=bf,skip=4pt}
\setlist{nosep,leftmargin=1.2em}

% IEEE-style section headings: Roman numerals, centered small caps for sections
\renewcommand{\thesection}{\Roman{section}}
\renewcommand{\thesubsection}{\Roman{section}-\Alph{subsection}}
\titlespacing*{\section}{0pt}{7pt}{3pt}
\titlespacing*{\subsection}{0pt}{5pt}{2pt}
\titleformat{\section}[block]{\normalfont\normalsize\bfseries\centering}{\thesection.}{0.5em}{\scshape}
\titleformat{\subsection}[block]{\normalfont\normalsize\bfseries}{\thesubsection.}{0.5em}{}

% Abstract styling
\renewcommand{\abstractnamefont}{\normalfont\bfseries}
\renewcommand{\abstracttextfont}{\normalfont\small\itshape}

\definecolor{codebg}{rgb}{0.96,0.96,0.96}
\definecolor{codecomment}{rgb}{0.35,0.45,0.35}
\definecolor{codekw}{rgb}{0.10,0.10,0.55}
\lstset{
  basicstyle=\ttfamily\scriptsize,
  backgroundcolor=\color{codebg},
  commentstyle=\color{codecomment}\itshape,
  keywordstyle=\color{codekw}\bfseries,
  language=C,
  breaklines=true,
  columns=fullflexible,
  keepspaces=true,
  frame=single,
  framesep=3pt,
  xleftmargin=2pt,
  aboveskip=4pt,
  belowskip=4pt,
  morekeywords={sqlite3,sqlite3_stmt,PyObject,size_t}
}

\setlength{\parskip}{2pt}
\setlength{\columnsep}{0.28in}
\pagestyle{empty}

\title{\vspace{-2.4em}\textbf{\large Pointer-Safe Marshalling at the Host--SQLite Boundary}\\[0.15em]
\normalsize\textnormal{Characterizing and Eliminating Lifetime-Contract Violations in Trusted Embedding Glue}}
\author{\normalsize Claude summary after chatting with the project lead}
\date{\vspace{-1.4em}}

\begin{document}
\maketitle
\thispagestyle{empty}

\begin{abstract}
\noindent SQLite is the most widely deployed database engine in the world, and
it runs as a library linked directly into its host: a fault in the engine is a
fault in the host's address space. Each host reaches the engine through
SQLite's C API. That boundary passes \emph{borrowed, short-lived pointers in
both directions}. Result values are lent out to the host, and buffers and
callbacks are lent in to the engine. The code that mediates this exchange is
trusted host code. A lifetime mistake in it, such as a borrowed pointer used
past a teardown or a callback context freed mid-invocation, is a memory-safety
vulnerability that no engine-level sandbox contains. We characterize the
boundary's API surface and show that the resulting defect class recurs across
nine independent host bindings spanning C, C++, Rust, and Go, while the
vulnerable boundary itself is replicated by every application that embeds
SQLite. We then argue that the obvious mitigation, copying every value across
the boundary, is too expensive for the engine's hottest path. We state the
goals for a marshalling layer that makes the contract violation
\emph{unrepresentable} while preserving zero-copy access, and argue that an
architectural capability model with linear capabilities and revocable
hierarchical delegation (Capstone) can realize such a layer.
\end{abstract}

\section{Introduction}

SQLite is embedded in browsers, mobile operating systems, language runtimes,
and countless applications; it ships as a C library compiled \emph{into} the
host process, with no process or memory isolation of its own. Consequently it
executes with the host's privileges in the host's address space, and any
memory-safety fault in or around the engine is a fault in the host. Throughout,
we use \emph{host} to mean any application or language binding that embeds
SQLite through its C API.

Every host reaches the engine across a single, narrow interface: the SQLite C
API. That interface is the subject of this paper, because the way data crosses
it is inherently lifetime-sensitive. Consider the most common operation a host
performs, reading a text column from a result row:

\begin{lstlisting}
sqlite3_step(stmt);                /* position on row 1 */
const unsigned char *s =
    sqlite3_column_text(stmt, 0);  /* borrow: ptr into engine memory */
sqlite3_step(stmt);                /* advance: row 1's buffer is freed */
puts((const char *)s);             /* use-after-free: s now dangles */
\end{lstlisting}

\noindent The pointer returned by \texttt{sqlite3\_column\_text} is not a copy:
it addresses a buffer \emph{owned by the engine}, and the contract states it
stays valid only until the next \texttt{step}, \texttt{reset}, or
\texttt{finalize} on that statement (or an implicit type conversion). A host
that keeps the pointer one step too long dereferences freed memory. Real
bindings have done exactly this by caching it for later use
(Section~\ref{sec:scope}). The hazard is symmetric: in the reverse direction,
buffers passed with \texttt{sqlite3\_bind\_*} and callbacks registered with the
engine are \emph{borrowed} across the call. They are owned by the host but held
by the engine.
Neither direction transfers ownership by default, and the contract is
documented but \emph{not enforced}: correctness rests entirely on each binding
author upholding it by hand. Table~\ref{tab:apis} enumerates the boundary's
principal API groups and the lifetime or ownership obligation attached to each.

\begin{table}[t]
\centering
\caption{Host--SQLite boundary API surface (the \texttt{sqlite3\_} prefix is
omitted). Each group crosses the boundary with a lifetime or ownership
obligation the host must uphold. Dir.: E$\rightarrow$H engine-to-host,
H$\rightarrow$E host-to-engine, cb callback.}
\label{tab:apis}
\scriptsize
\renewcommand{\arraystretch}{1.25}
\begin{tabular}{@{}p{2.15cm} c p{3.45cm}@{}}
\toprule
\textbf{API group} & \textbf{Dir.} & \textbf{Boundary concern} \\
\midrule
\texttt{column\_blob/text} & E$\rightarrow$H & Borrowed ptr; invalid after \texttt{step}/\texttt{reset}/\texttt{finalize} or type conversion. \\
\texttt{column\_value} & E$\rightarrow$H & \emph{Unprotected} value; unsafe except via \texttt{bind/result\_value}. \\
\texttt{column\_name} & E$\rightarrow$H & Borrowed string; invalid on reprepare/finalize. \\
\texttt{value\_blob/text} & E$\rightarrow$H & Valid only within the UDF invocation. \\
\texttt{bind\_blob/text} & H$\rightarrow$E & Ownership flag: \texttt{STATIC} (caller owns), \texttt{TRANSIENT} (copy), or destructor. \\
\texttt{result\_*} (UDF) & H$\rightarrow$E & Same ownership rules on the return path. \\
\texttt{\{bind,result,}\newline\texttt{value\}\_pointer} & H$\leftrightarrow$E & Type-tagged raw pointer threaded through SQL. \\
\texttt{create\_function},\newline hooks, \texttt{set\_}\newline\texttt{authorizer}, collation, progress, \texttt{trace\_v2} & cb & Host code invoked later; context + destructor lifetime. \\
\texttt{blob\_open/}\newline\texttt{read/write} & E$\leftrightarrow$H & Incremental handle invalidated by writes to the row. \\
\texttt{\{de\}serialize} & H$\rightarrow$E & Hands engine a memory image as a DB; \texttt{FREEONCLOSE} ownership. \\
\texttt{prepare}/\texttt{step}/\newline\texttt{finalize}, \texttt{close\_v2} & H$\leftrightarrow$E & Statement/handle lifecycle; ``zombie'' connection on close with live stmts. \\
\bottomrule
\end{tabular}
\end{table}

This boundary is not a niche concern: it is replicated, independently, by every
host that embeds SQLite. Table~\ref{tab:hosts} lists twelve widely deployed
hosts and the binding through which each crosses the API. The bindings are
written
in C, C++, Rust, Go, and others, but they all confront the same lifetime
contract, and each re-implements the marshalling glue that enforces it.

\begin{table}[t]
\centering
\caption{Representative hosts embedding SQLite and their bindings. The
marshalling glue is re-implemented per binding.}
\label{tab:hosts}
\footnotesize
\renewcommand{\arraystretch}{1.15}
\begin{tabular}{@{}p{2.7cm} p{3.0cm} l@{}}
\toprule
\textbf{Host} & \textbf{Binding} & \textbf{Lang} \\
\midrule
Chrome / Chromium   & \texttt{sql::Database} (forked) & C++ \\
Safari / WebKit     & WebSQL / storage layer          & C++ \\
Mozilla Firefox     & mozStorage                      & C++ \\
Android (AOSP)      & \texttt{android\_database\_}\allowbreak\texttt{SQLiteConnection} (JNI) & C++/Java \\
CPython             & \texttt{\_sqlite3} module        & C \\
PHP                 & \texttt{ext/sqlite3}, \texttt{pdo\_sqlite} & C \\
Ruby                & \texttt{sqlite3-ruby} gem        & C \\
Go                  & \texttt{mattn/go-sqlite3} (cgo)  & C/Go \\
Rust                & \texttt{rusqlite}, \texttt{diesel} & Rust \\
Node.js             & \texttt{better-sqlite3}, \texttt{node-sqlite3} & C++ \\
Signal              & SQLCipher (forked)               & C \\
.NET                & \texttt{SQLitePCLRaw}            & C\#/C \\
\bottomrule
\end{tabular}
\end{table}

\subsection{Threat model}
The danger is structural and independent of any sandboxing strategy. The
marshalling glue of Table~\ref{tab:hosts} is trusted host code: it runs with
full host privilege and mediates every value entering or leaving the engine.
Consider a lifetime error in that glue: a result pointer dereferenced after the
statement that owned it advanced, or a callback context freed while the engine
still holds it. This is a use-after-free or related memory-safety vulnerability.
It is reachable whenever an adversary influences the data or control flow that
crosses the boundary. Because the contract is documented but not enforced,
correctness rests entirely on each binding author manually upholding it, and
the defect recurs (Section~\ref{sec:scope}).

\subsection{Memory-safe isolation is not a complete answer}
A natural response is to isolate the engine itself. One can run it in a
separate process, or compile it to WebAssembly, so that malicious input
corrupts only the engine's own sandbox. This is valuable, but it does not
solve the problem stated above. Engine-level isolation contains faults
\emph{inside} the engine. The marshalling glue sits on the trusted host side
of the boundary, outside any such sandbox. A use-after-free in the host's
result-reading or callback-dispatch code is unaffected by where the engine
runs. Worse, the safe way to traverse a hard isolation boundary is to serialize
and copy every value across it, which imposes exactly the overhead we argue
against in Section~\ref{sec:goals}. Isolation therefore relocates the trusted
boundary. It does not eliminate the lifetime contract that this paper targets.

\section{Problem Scope: A C Host Example}
\label{sec:scope}

The defect class is concrete and cross-ecosystem, and it is not limited to the
read-side hazard of Section~I. A representative recent instance on the
\emph{callback} path is CPython's \texttt{\_sqlite3} module (gh-142830, 2025), a
C extension wrapping the SQLite C API. A progress-handler callback that returns
an object with a custom \texttt{\_\_bool\_\_} can de-register itself
\emph{while the callback is still executing}, freeing the callback context the
engine then dereferences:

\begin{lstlisting}
/* host callback context freed mid-invocation */
ret = PyObject_CallNoArgs(callable);
rc  = PyObject_IsTrue(ret);   /* runs __bool__:
        set_progress_handler(None) -> frees ctx */
if (rc < 0)
    print_or_clear_traceback(ctx); /* UAF: ctx freed */
\end{lstlisting}

\noindent AddressSanitizer reports a heap-use-after-free read inside the host's
callback dispatch, reached directly through \texttt{sqlite3VdbeExec}. The flaw
is in the \emph{binding's} memory management, the trusted-glue surface of
Section~I, not in SQLite.

Table~\ref{tab:scope} lists 19 real, publicly tracked boundary-contract defects
across host ecosystems. Two root causes dominate: \textbf{callbacks/hooks}
registered into the engine, and \textbf{statements or connections used across
teardown}. Conspicuously rare is the plain read-out-a-column path, the part a
copy or a bounds-checked read already makes safe. The pattern defines the
problem statement: \emph{the trusted marshalling layer is a recurring,
exploitable memory-safety surface that engine isolation leaves exposed.}

\begin{table*}[t]
\centering
\caption{Host-side SQLite boundary-contract defects, across ecosystems.
UAF\,=\,use-after-free; UAC\,=\,use-after-close. Rows 1--11 are
UAF, double-free, use-after-close, or finalize-time memory crashes; rows
12--16 are adjacent lifetime, callback, or concurrency-safety defects; rows
17--19 are bind/result/lifecycle mishandling (crash or stale-state) in the
host glue. All links verified.}
\label{tab:scope}
\scriptsize
\renewcommand{\arraystretch}{1.18}
\begin{tabular}{@{}r p{2.0cm} p{2.0cm} p{4.0cm} l@{}}
\toprule
\textbf{\#} & \textbf{Host (lang)} & \textbf{Ref} & \textbf{Class} & \textbf{Link} \\
\midrule
1  & CPython (C)        & gh-142830                 & UAF, progress-handler context freed mid-call & \href{https://github.com/python/cpython/issues/142830}{\nolinkurl{github.com/python/cpython/issues/142830}} \\
2  & rusqlite (Rust)    & RUSTSEC-2021-0128         & UAF, closure lifetimes on functions/hooks    & \href{https://rustsec.org/advisories/RUSTSEC-2021-0128.html}{\nolinkurl{rustsec.org/advisories/RUSTSEC-2021-0128.html}} \\
3  & diesel (Rust)      & RUSTSEC-2021-0037         & UAF, column-name pointer cached across step  & \href{https://rustsec.org/advisories/RUSTSEC-2021-0037.html}{\nolinkurl{rustsec.org/advisories/RUSTSEC-2021-0037.html}} \\
4  & PHP (C)            & Bug \#66550               & UAC, statement used after \texttt{sqlite3\_close} & \href{https://bugs.php.net/bug.php?id=66550}{\nolinkurl{bugs.php.net/bug.php?id=66550}} \\
5  & PHP (C)            & Bug \#69971               & UAF, statement/database destruction order    & \href{https://bugs.php.net/bug.php?id=69971}{\nolinkurl{bugs.php.net/bug.php?id=69971}} \\
6  & PHP (C)            & Bug \#77977               & UAF, via user-defined function               & \href{https://bugs.php.net/bug.php?id=77977}{\nolinkurl{bugs.php.net/bug.php?id=77977}} \\
7  & CPython (C)        & gh-99886                  & UAF, subclassed Cursor dealloc order         & \href{https://github.com/python/cpython/issues/99886}{\nolinkurl{github.com/python/cpython/issues/99886}} \\
8  & CPython (C)        & gh-85981 / bpo-41815      & UAC, backup() on a closed connection         & \href{https://github.com/python/cpython/issues/85981}{\nolinkurl{github.com/python/cpython/issues/85981}} \\
9  & sqlite3-ruby (C)   & Issue \#49                & Finalize-time segfault on close              & \href{https://github.com/sparklemotion/sqlite3-ruby/issues/49}{\nolinkurl{github.com/sparklemotion/sqlite3-ruby/issues/49}} \\
10 & sqlite3-ruby (C)   & "Closed Statement"        & UAC, statement reused after close            & \href{https://groups.google.com/g/sqlite3-ruby/c/SGRQE_2MZ8I}{\nolinkurl{groups.google.com/g/sqlite3-ruby/c/SGRQE_2MZ8I}} \\
11 & go-sqlite3 (cgo)   & PR \#1303                 & Double-free / finalizer not removed on Close & \href{https://github.com/mattn/go-sqlite3/pull/1303}{\nolinkurl{github.com/mattn/go-sqlite3/pull/1303}} \\
12 & expo-sqlite (C/JNI)& PR \#34992                & Unfinalized statements $\rightarrow$ NPE on close & \href{https://github.com/expo/expo/pull/34992}{\nolinkurl{github.com/expo/expo/pull/34992}} \\
13 & CPython (C)        & gh-149738                 & Null-deref, deleted \texttt{row\_factory}     & \href{https://github.com/python/cpython/issues/149738}{\nolinkurl{github.com/python/cpython/issues/149738}} \\
14 & CPython (C)        & bpo-31746 / PR \#27472    & Null-deref, uninitialised Connection         & \href{https://github.com/python/cpython/pull/27472}{\nolinkurl{github.com/python/cpython/pull/27472}} \\
15 & CPython (C)        & gh-101767                 & Callback GIL-state unsafe in subinterpreters & \href{https://github.com/python/cpython/issues/101767}{\nolinkurl{github.com/python/cpython/issues/101767}} \\
16 & datasette-authz (C)& Issue \#3                 & Authorizer callback-context lifetime churn   & \href{https://github.com/datasette/datasette-sqlite-authorizer/issues/3}{\nolinkurl{github.com/datasette/datasette-sqlite-authorizer/issues/3}} \\
17 & node-sqlite3 (C++) & Issue \#1449              & Fatal abort, BindParameter typecast failure  & \href{https://github.com/TryGhost/node-sqlite3/issues/1449}{\nolinkurl{github.com/TryGhost/node-sqlite3/issues/1449}} \\
18 & PHP (C)            & Bug \#79294               & Stale-state read, columnType after reset     & \href{https://bugs.php.net/bug.php?id=79294}{\nolinkurl{bugs.php.net/bug.php?id=79294}} \\
19 & PHP (C)            & PR \#5204 / Bug \#64531   & Result/statement lifecycle mismatch          & \href{https://github.com/php/php-src/pull/5204}{\nolinkurl{github.com/php/php-src/pull/5204}} \\
\bottomrule
\end{tabular}
\end{table*}

\section{Solution Goals}
\label{sec:goals}

A solution should make the host--SQLite boundary safe \emph{by construction},
so that no lifetime mistake in hand-written glue can become a host compromise.
We state two goals.

\subsection{G1: Protect pointers crossing in both directions}
The boundary must enforce, rather than document, the ownership and lifetime
contract of Table~\ref{tab:apis}. Engine$\rightarrow$host: a value handed out
must remain valid for a well-defined scope and be impossible to dereference
after the owning statement advances or is torn down. Host$\rightarrow$engine:
bound buffers and registered callbacks must have their lifetimes pinned to the
engine's use of them, so that a callback context cannot be freed while the
engine still holds it (the gh-142830 pattern), and a statement cannot be used
after its connection closes (the PHP and Ruby patterns). Enforcement should be
mechanical: scoped or owned handles, lifetime typing, or a representation in
which a dangling reference cannot be named.

\subsection{G2: Preserve performance: copying is not the answer}
The obvious safe default is to copy every value across the boundary, using
\texttt{SQLITE\_TRANSIENT} semantics or serialize-and-deserialize across a
process or WebAssembly boundary. That trades safety for throughput, and the
trade is too expensive on the engine's hottest path. Every row read and every
parameter bound incurs an allocation and a memory copy. Large \texttt{BLOB} and
\texttt{TEXT} columns copy in full, doubling memory traffic and defeating
SQLite's zero-copy read design. A solution must instead protect the borrowed
pointer \emph{in place}, through enforced lifetimes, scoped or owned handles, or
a memory-safe boundary representation, so that the common read and bind path
stays copy-free while the dangling dereference becomes unrepresentable.

\medskip
\noindent\textbf{Success criterion.} A marshalling layer in which the
boundary-contract defects of Table~\ref{tab:scope} cannot be written: borrowed
pointers are lifetime-checked in both directions (G1), and the zero-copy fast
path is retained for the common read and bind operations (G2).

\section{Proposed Solution: A Capability-Based\\Marshalling Layer}
\label{sec:soln}

We sketch how Capstone can enforce the contract of Table~\ref{tab:apis} by
construction, meeting G1 and G2 without copying. Capstone is an implemented
capability ISA with \emph{linear capabilities}, explicit \emph{revocation
capabilities}, and \emph{uninitialised capabilities}~\cite{capstone,capstonespec}.
We model the engine and the host binding as two cooperating domains \emph{in one
address space}. Every value crossing the boundary is named by a capability
handle to the real buffer, not a raw pointer. Capstone's stated motivation for
revocable delegation is exactly our scenario: after an untrusted component is
granted access to a region, the grantor must be able to revoke
it~\cite{capstonespec}. Capabilities fit because they are \emph{trustless}.
Neither domain need trust the other's lifetime discipline. This is the situation
that defeats engine-level isolation (Section~I-B), where the glue is trusted but
buggy and copying across a hard boundary is the cost G2 forbids.

\subsection{Linear capabilities replace borrowed pointers}
A linear capability is \emph{alias-free} and move-only. Instructions may move
but not copy it, so while one domain holds it no other capability grants access
to the region~\cite{capstonespec}. Replace the borrowed result pointer of
Section~I with a linear capability that the engine delegates to the host for
the borrow window. Three consequences follow. (i)~The dangling read becomes
\emph{unrepresentable}: access is only via the handle, so a host that
over-retains it cannot silently read freed memory. (ii)~Double-free becomes
structurally impossible. Double-free and double-finalize (rows~9,~11) are
aliasing bugs that require two live references. A move-only linear capability
admits exactly one holder, so a second \texttt{finalize} has nothing to consume.
(iii)~Zero-copy is preserved (G2): the capability addresses the engine's actual
buffer. The reverse direction is symmetric. The host delegates a linear
capability for a bound buffer, and exclusivity forbids it from freeing the
buffer out from under the engine. This removes the very reason
\texttt{TRANSIENT} copies exist.

\subsection{Revocation capabilities end the borrow}
Capstone's revocation is explicit and is the natural model for the contract's
end-of-life points. A revocation capability carries no access of its own; it is
minted from a linear capability and later used to invalidate the capabilities
derived from it~\cite{capstonespec}. Before lending a result buffer, the engine
mints a revocation capability for it and retains it, delegating the linear
capability to the host. At each point the contract says the borrow ends, the
engine \emph{executes a revoke}. Those points are the next \texttt{step},
\texttt{reset}, or \texttt{finalize} on that statement, or an in-place type
conversion. The revoke sets the host's borrowed capability invalid
(\texttt{valid}\,=\,0). A subsequent dereference then \emph{faults} rather than
reading freed or rewritten memory (rows~1--8,~13,~14). This places no new burden
on the engine. It already runs code at exactly these transitions, so minting at
hand-out and revoking at teardown fall on paths it already executes. The reverse
direction reuses the same primitive together with \emph{uninitialised
capabilities}. A buffer the engine reclaims from the host is returned as an
uninitialised, write-only capability that becomes usable only once fully
overwritten, so the reclaiming side cannot read the previous occupant's
data~\cite{capstonespec}. Callback dispatch is a domain crossing, where the
engine invokes host code holding a context. Capstone provides sealed and
sealed-return capabilities for safe context switching here. The context's
capabilities are preserved across the call, and a revoke of the context (the
gh-142830 re-entrant free) invalidates the engine's handle to it. We stress
this is \emph{safe-fail}: it converts a use-after-free into a deterministic
capability fault, eliminating the exploitable vulnerability rather than making
the buggy program complete.

\subsection{Hierarchical delegation contains the glue}
The boundary forms an ownership tree (connection $\supset$ statement $\supset$
value), and bindings often re-delegate borrowed values onward to application
code. Capstone realises this directly. Revocation capabilities may overlap, and
among aliasing ones it defines a seniority by creation order. The earlier one
($c \prec_t d$) is the stronger, and revoking with it invalidates the weaker,
overlapping ones. This gives arbitrarily extensible chains of
delegation~\cite{capstonespec}. The engine mints a senior revocation capability
covering a statement's region when it is prepared. \texttt{close}, executed as
a revoke with that senior capability, invalidates every outstanding statement
and value capability beneath it. This resolves the use-after-close class
(rows~4,~8,~12) by construction. More fundamentally, the engine's senior
revocation \emph{dominates} the binding's. A buggy binding that re-delegates a
borrowed value and forgets to revoke it is still cleaned up by the engine's
revoke at the next lifecycle transition. Enforcement thus moves off ``every
binding author upholds the contract by hand,'' the threat model's core failure,
and onto an invariant the buggy glue cannot defeat.

\subsection{Applying the discipline to the observed defects}
Table~\ref{tab:fix} maps every defect of Table~\ref{tab:scope} to the capability
mechanism that makes it safe. Most rows are direct. A borrowed handle is revoked
at a defined point, so a later use faults instead of corrupting memory. Three
cases deserve comment. Row~14 is use-before-initialisation rather than a borrow,
so it uses an uninitialised capability that carries no read authority until
setup completes. Rows~18 and~19 are stale-state reads. They go through only once
the borrow discipline is extended to \emph{derived} state, so that column
metadata and result objects are themselves row- or statement-scoped borrows that
the engine revokes on \texttt{reset} or re-\texttt{step}. Row~17 is the weakest
fit: capabilities do not repair the faulty type coercion, but a bounded
read-only borrow of the parameter turns a fatal process abort into a local,
trappable fault. Only row~15 is out of reach. A subinterpreter GIL deadlock is a
liveness defect with no pointer lifetime.

\begin{table*}[t]
\centering
\caption{Capability-based fix for each defect of Table~\ref{tab:scope}.
Primitives: \textbf{L} linear (exclusive, move-only borrow); \textbf{R}
revocation (explicit revoke at the borrow's end); \textbf{H} hierarchical revoke
(a senior revoke cascades on \texttt{close} or destroy); \textbf{U} uninitialised
(no use before initialisation); \textbf{S} sealed (safe callback domain switch).
All defects but row~15 are addressed.}
\label{tab:fix}
\footnotesize
\renewcommand{\arraystretch}{1.15}
\setlength{\tabcolsep}{5pt}
\begin{tabular}{@{}r l p{14.2cm}@{}}
\toprule
\textbf{\#} & \textbf{Prim.} & \textbf{Capability-based fix} \\
\midrule
1  & L,R,S & The engine borrows the callback context for the call. A mid-call free revokes it, so the engine's later use faults. \\
2  & L,R,S & The closure is lent while registered. Dropping its captures revokes the handle, so a later invocation faults. \\
3  & L,R   & The column pointer is a row-scoped borrow. \texttt{step} revokes it, so a cached read faults. \\
4  & H     & The statement is a capability under the connection. \texttt{close} revokes the subtree, so reuse faults. \\
5  & H     & The statement handle is derived from the database. Destroying the database revokes it, so neither destruction order leaves a dangling handle. \\
6  & L,R,S & The function and its data are lent while registered. Freeing revokes them, so the engine's call faults. \\
7  & H     & The cursor holds derived connection and statement handles. Owner teardown revokes them, so out-of-order dealloc faults. \\
8  & R     & \texttt{backup} borrows the connection. \texttt{close} revokes the borrow, so the call faults. \\
9  & H     & The statement is a capability under the connection. \texttt{close} revokes it, so \texttt{finalize} faults instead of corrupting memory. \\
10 & H     & As row~4: a reuse after \texttt{close} hits an invalid capability. \\
11 & L     & The handle is one move-only capability. \texttt{Close} consumes it, so no second alias remains to free again. \\
12 & H     & Statements are capabilities under the connection. \texttt{close} revokes the subtree, so no unfinalised state can block the close. \\
13 & L,R   & The cursor borrows the \texttt{row\_factory} handle. Deleting it revokes the borrow, so the later use faults. \\
14 & U     & The connection is an uninitialised capability until setup completes. A premature access has no read authority and faults. \\
15 & n/a   & Not addressed. A subinterpreter GIL deadlock is a liveness defect with no pointer lifetime. \\
16 & L,R,S & The authorizer context is lent while registered. Mint at register and revoke at unregister fix the lifetime across language versions. \\
17 & L     & The bound value crosses as a bounded read-only borrow. A malformed access faults locally instead of aborting the process. \\
18 & R     & The column type is a row-scoped borrow. \texttt{reset} revokes it, so a later read faults instead of returning stale data. \\
19 & L,R   & The result is a borrow derived from the statement. \texttt{reset} or re-stepping revokes it, so a mismatched fetch faults. \\
\bottomrule
\end{tabular}
\end{table*}

\subsection{Scope and limitations}
The model meets G1 (both directions capability-mediated, stale handles trap,
double-free impossible, \texttt{close} cascades) and G2 (handles, not copies,
and the safe borrow replaces the mandatory \texttt{TRANSIENT} copy). Capstone
overhead is reported below 50\% in common cases~\cite{capstone}. Two honest
costs remain. First, linear capabilities are exclusive. Patterns that genuinely
require shared concurrent reads of one value must fall back to non-linear
capabilities, which forfeit the alias-freedom on which the clean revocation
guarantee rests. The borrow-then-reclaim pattern that dominates
Table~\ref{tab:scope} does not need such sharing. Second, revocation here is
explicit. The engine must mint and revoke at the lifecycle transitions. It
already executes code there, but the obligation is real. A capability design
that made revocation implicit on the owner's next access, as the unimplemented
CapsLock proposes for mixed Rust and FFI~\cite{capslock}, would discharge that
obligation automatically. We note it as future work rather than rely on it.

\section{Conclusion}
SQLite's safety inside an embedding host turns on a boundary the engine does not
control: the marshalling glue that lends pointers in both directions. We
characterized that boundary, showed its lifetime contract is documented but not
enforced, and collected 19 real defects across nine bindings that share one
shape. A borrowed handle is used past the window in which it is valid.
Engine-level isolation does not address this, because the glue is trusted and
sits outside any sandbox, and copying every value defeats SQLite's zero-copy
design. We argued that Capstone can enforce the contract in place. Linear
capabilities make a borrow exclusive and a double-free unrepresentable. Explicit
revocation ends the borrow at the points the contract already names.
Senior revocation makes a connection's \texttt{close} cascade to its statements.
Uninitialised capabilities cover safe reclaim and use-before-init. Table~\ref{tab:fix}
maps all but one of the defects to one of these mechanisms; the exception is a
deadlock, which no memory model addresses. The outcome is \emph{safe-fail}: a
contract violation becomes a deterministic fault rather than an exploitable
corruption, with the zero-copy fast path preserved. A full result rests on two
deliverables we leave to evaluation: a binding co-designed with the engine so
the engine mints and revokes the capabilities, and a measurement showing the
capability-mediated borrow path stays close to raw pointers and well below the
copy baseline.

\begin{thebibliography}{9}
\footnotesize
\bibitem{capstone} J.~Z.~Yu, C.~Watt, A.~Badole, T.~E.~Carlson, and P.~Saxena.
``Capstone: A Capability-based Foundation for Trustless Secure Memory Access.''
\emph{USENIX Security}, 2023.
\href{https://www.usenix.org/conference/usenixsecurity23/presentation/yu-jason}{\nolinkurl{usenix.org/conference/usenixsecurity23/presentation/yu-jason}}
\bibitem{capstonespec} ``The Capstone-RISC-V Instruction Set Reference,'' v1.0.
\href{https://capstone.kisp-lab.org/specs/}{\nolinkurl{capstone.kisp-lab.org/specs}}
\bibitem{capslock} J.~Z.~Yu, C.~Watt, et al.
``Securing Mixed Rust with Hardware Capabilities'' (CapsLock).
\emph{ACM CCS}, 2025. Related/future work; not relied upon here.
\href{https://arxiv.org/abs/2507.03344}{\nolinkurl{arxiv.org/abs/2507.03344}}
\end{thebibliography}

\end{document}
